\chapter{Introduction}
\label{ch:introduction}

\section{Background and motivation}

Advances in neural text-to-speech (TTS), singing-voice synthesis, text-to-music systems, and general audio generative models have lowered the effort required to produce convincing synthetic audio. Such clips can be shared at scale through social media, messaging applications, and voice-based interfaces, where listeners often lack tools to verify provenance. \textbf{Audio deepfake detection (ADD)} is the learning problem of estimating whether a recording is \emph{bona fide} (authentic) or \emph{synthetic/spoof} (fake), usually from acoustic features alone.

Early community effort focused on attacks against automatic speaker verification, with shared datasets and standardized metrics such as the equal error rate (EER)~\cite{kinnunen2017detection,wang2020asvspoof}. More recent benchmarks explicitly target \textbf{all-type} content (speech, singing, music, and environmental sound) and realistic channel perturbations~\cite{xie2026atadd}. This project adopts AT-ADD \textbf{Track~2}, in which the detector must classify real versus fake without knowing the audio category at test time. This setting is closer to open-world deployment than speech-only laboratory evaluation, but it also increases the risk that a model tuned on one domain fails on another.

Deployment adds a second motivation. Audio that reaches a detector has usually been re-encoded, resampled, or mixed with noise, and the detector is often expected to run on phones, browsers, or embedded devices with limited memory and energy. Self-supervised encoders with hundreds of millions of parameters dominate the top AT-ADD entries~\cite{xie2026atadd}, but they are difficult to run on such hardware. These observations motivate a compact detector whose robustness to channel changes is obtained during training, so that the inference cost stays that of a small network.

\section{Challenges in audio deepfake detection}

Building a practical audio deepfake detector is difficult for several interacting reasons rather than because of binary classification alone. The main challenges addressed in this work are the following.

\paragraph{Diverse acoustic categories:} Speech, singing, music, and environmental sound differ in harmonic structure, temporal dynamics, and signal-to-noise characteristics. Artifacts that indicate synthesis in speech (for example, phase inconsistencies in vocoded speech) may not transfer to music generated by diffusion or transformer-based models. A detector that exploits speech-specific cues can perform well on one audio type yet fail on music or environmental sound, so detection performance must be examined for each audio type rather than only in aggregate.

\paragraph{Generator diversity:} Fake clips may originate from autoregressive TTS, flow-based vocoders, singing conversion, text-to-music, or general audio synthesis. Protocols that hold out generators at test time measure generalization to unseen synthesis pipelines~\cite{xie2026atadd}; a model may otherwise memorize generator-specific fingerprints rather than learn invariant signs of synthesis.

\paragraph{Robustness under realistic audio conditions:} Benchmark clips are often stored in relatively high-quality formats, whereas deployed audio is frequently re-encoded, resampled for streaming, mixed with background noise, or filtered by telephone channels. Features that are useful on clean spectrograms may be attenuated or aliased after such processing, producing a \textbf{domain shift} between training assumptions and application conditions. Claims of compression robustness therefore require evaluation under explicit degradations, not clean scores alone.

\paragraph{Spurious cues:} Classifiers can reach high accuracy by exploiting recording conditions, codecs, or dataset biases instead of genuine evidence of synthesis. Such shortcuts tend to break when the channel changes, so robustness must be evaluated explicitly rather than inferred from clean validation scores.

\paragraph{Computational constraints:} Mobile, embedded, and browser-side inference limit parameter count, memory bandwidth, and energy per clip. Compactness is a necessary design constraint for edge use, but it is not sufficient: real-time suitability must be confirmed by measurement on a target platform.

\section{Problem statement}

Given a mono audio segment of fixed duration (four seconds at 24\,kHz), estimate the label $y \in \{0,1\}$ (real $=0$, fake $=1$). AT-ADD Track~2 mixes four audio types in evaluation, and the type is not provided at inference. The problem is to design a detector that (i) works across all four audio types, (ii) remains stable when the input has been degraded by compression, noise, resampling, or band-limiting, and (iii) keeps an inference cost compatible with edge deployment.

The \textbf{primary} experimental objective is clean held-out evaluation on AT-ADD Track~2 under a fixed, reproducible preprocessing and training protocol.

\section{Objectives}

The objectives of this project are to:
\begin{enumerate}
  \item implement and validate a reproducible acoustic preprocessing pipeline and lightweight baseline architectures (pure 2D CNN vs.\ recurrent CRNNs) on speech (ASVspoof 2019) and multi-domain audio (AT-ADD Track~2), benchmarking against published literature;
  \item design and implement Frequency-Time Robustness Coherence Learning (FT-RCL), which uses paired clean and degraded views, prediction coherence, band-wise feature signature coherence, a severity curriculum, and gated robustness weights, without changing the inference graph;
  \item implement the STAR-Edge extension of FT-RCL, which adds asymmetric bona-fide robustness, severity trajectory coherence, and channel auxiliary detection;
  \item evaluate all systems with Macro-F1, accuracy, AUROC, and EER on clean held-out audio, overall and per audio type, and log full training histories;
  \item document checkpoint size and the planned edge-deployment measurements (latency, memory, INT8 accuracy).
\end{enumerate}

\section{Contributions}

\paragraph{Methodological contributions:}
\begin{itemize}
  \item A unified coherence-based robustness training framework that jointly applies dual-view classification, prediction coherence, band-wise signature coherence, severity scheduling, and batch-mean robustness gating to a shared lightweight backbone, leaving the inference-time architecture identical to that of the baseline.
  \item A separation of waveform-domain degradation from the mel frontend parameters, which allows a controlled comparison between clean-only training and coherence-based training.
  \item A robustness-aware model selection criterion that balances clean development performance against the estimated metric drop under degradation.
  \item The STAR-Edge extension, which augments FT-RCL with asymmetric bona-fide robustness, severity trajectory coherence, and a channel auxiliary detection branch.
\end{itemize}

\clearpage
\paragraph{Experimental contributions:}
\begin{itemize}
  \item A foundational baseline investigation evaluating pure 2D CNN versus recurrent CRNN architectures on full-scale multi-domain audio, identifying duration shortcut artifacts on ASVspoof 2019, and establishing competitive performance against published literature (Pham et al., 2024).
  \item A clean-evaluation comparison of the baseline and FT-RCL on AT-ADD Track~2, documenting aggregate and per-type trade-offs, including the cases where FT-RCL does not outperform the baseline.
  \item Documented training and clean-evaluation protocols, with ablation and edge-benchmarking directions for follow-on work.
\end{itemize}
